Work under the witness constructed in Propositions~\(\mathrm{prop:endogenous\text{-}qv\text{-}self\text{-}normalization\text{-}obstruction}\) and~\(\mathrm{prop:saco\text{-}theorem\text{-}5\text{-}14\text{-}counterexample}\), and take the even horizons \(T=2n\).  Write \(\mathcal F_t\) for the witness filtration.  During the first block the construction gives independent scores
\[
 P(\xi_t=2)=P(\xi_t=-2)=\frac12,
 \qquad 1\le t\le n,
\tag{1}
\]
and the sign of \(S_n=\sum_{t=1}^n\xi_t\) selects the second-block conditional variance: it is \(16/3\) when \(S_n>0\) and \(4\) when \(S_n\le0\).  The first-block conditional variance is \(4\).  These properties are unchanged by the bounded AIPW augmentation because Proposition~\(\mathrm{prop:saco\text{-}theorem\text{-}5\text{-}14\text{-}counterexample}\) identifies its centered AIPW increment pathwise with \(\xi_t\).

Let \(k=k_{2n}\), where \(k_T=o(T)\) is deterministic.  Then \(k/n\to0\), so \(k<n\) for all sufficiently large \(n\).  Define the terminal regime label and its conditional success probability by
\[
 J_n=\mathbf 1\{S_n>0\},
 \qquad
 p_n=P(J_n=1\mid\mathcal F_k).
\tag{2}
\]
We first show that no \(\mathcal F_k\)-measurable binary variable can predict \(J_n\) consistently.  Conditional on \(\mathcal F_k\), equation~(1) gives
\[
 S_n=S_k+2R_{n-k},
 \qquad
 R_r=\sum_{i=1}^r\epsilon_i,
\tag{3}
\]
where the \(\epsilon_i\) are conditionally independent symmetric signs.  The maximum point mass of \(R_r\) is bounded by \(C r^{-1/2}\) for a universal finite \(C\).  To verify rather than assume this bound, if \(r=2\ell\), then
\[
 \max_xP(R_{2\ell}=x)
 =2^{-2\ell}\binom{2\ell}{\ell}
 =\prod_{j=1}^{\ell}\frac{2j-1}{2j}
 \le\frac1{\sqrt{\ell+1}}.
\tag{4}
\]
The last inequality follows by induction, because
\[
 \frac{2\ell+1}{2\ell+2}
 \le \sqrt{\frac{\ell+1}{\ell+2}}.
\]
For odd \(r=2\ell+1\), either modal probability is \((2\ell+1)/(2\ell+2)\) times the modal probability at \(2\ell\), so the same \(Cr^{-1/2}\) bound follows after enlarging \(C\).

By symmetry, \(|P(R_r>0)-1/2|\) is at most one modal mass.  Moving the threshold in \(P(2R_{n-k}>-S_k\mid\mathcal F_k)\) from zero to \(-S_k\) crosses at most \(1+|S_k|\) lattice atoms.  Hence (3)--(4) imply
\[
 \left|p_n-\frac12\right|
 \le \frac{C(1+|S_k|)}{\sqrt{n-k}}.
\tag{5}
\]
For \(s<t\le k\), the martingale-difference property in the witness gives
\[
 E(\xi_s\xi_t)
 =E\{\xi_sE(\xi_t\mid\mathcal F_{t-1})\}=0.
\]
Together with (1), this yields
\[
 E(S_k^2)=\sum_{t=1}^kE(\xi_t^2)=4k,
 \qquad
 E|S_k|\le\sqrt{E(S_k^2)}=2\sqrt{k}.
\tag{6}
\]
Taking expectations in (5) and using \(k/n\to0\) therefore gives
\[
 E\left|p_n-\frac12\right|
 \le \frac{C(1+2\sqrt{k})}{\sqrt{n-k}}
 \longrightarrow0.
\tag{7}
\]
If \(B_n\) is any \(\mathcal F_k\)-measurable \(\{0,1\}\)-valued predictor, conditioning on \(\mathcal F_k\) gives
\[
 P(B_n\ne J_n\mid\mathcal F_k)
 =\mathbf1\{B_n=0\}p_n+
   \mathbf1\{B_n=1\}(1-p_n)
 \ge \min\{p_n,1-p_n\}.
\]
Consequently, by (7),
\[
 P(B_n\ne J_n)
 \ge E\min\{p_n,1-p_n\}
 =\frac12-E\left|p_n-\frac12\right|
 \longrightarrow\frac12.
\tag{8}
\]

We now calculate the suffix predictable quadratic variation exactly.  There are \(n-k\) retained first-block increments, each with conditional variance \(4\), and \(n\) second-block increments, each with conditional variance \(16/3\) on \(\{J_n=1\}\) and \(4\) on \(\{J_n=0\}\).  Thus
\[
 q_n:=\frac1{2n}\sum_{t=k+1}^{2n}
 E(\xi_t^2\mid\mathcal F_{t-1})
 =
 \begin{cases}
 \displaystyle\frac{14}{3}-\frac{2k}{n},&J_n=1,\\[5pt]
 \displaystyle4-\frac{2k}{n},&J_n=0.
 \end{cases}
\tag{9}
\]
The two possible values in (9) differ by \(2/3\), and \(0<q_n<5\) for all sufficiently large \(n\).

Suppose, for contradiction, that there were positive \(\mathcal F_k\)-measurable random variables \(\Lambda_{2n}\) such that the stabilization displayed in the proposition held.  By the definition of \(q_n\), that stabilization is precisely
\[
 \frac{q_n}{\Lambda_{2n}}\longrightarrow_P1.
\tag{10}
\]
On the event \(\{|q_n/\Lambda_{2n}-1|<1/2\}\), positivity and \(q_n<5\) imply \(\Lambda_{2n}<2q_n<10\).  On that event,
\[
 |\Lambda_{2n}-q_n|
 =\Lambda_{2n}\left|\frac{q_n}{\Lambda_{2n}}-1\right|
 \le10\left|\frac{q_n}{\Lambda_{2n}}-1\right|.
\tag{11}
\]
Equation~(10) makes the event in (11) have probability tending to one, so (11) proves
\[
 |\Lambda_{2n}-q_n|\longrightarrow_P0.
\tag{12}
\]
Define the \(\mathcal F_k\)-measurable binary predictor
\[
 B_n=\mathbf1\left\{\Lambda_{2n}>
 \frac{13}{3}-\frac{2k}{n}\right\}.
\tag{13}
\]
The threshold in (13) is the midpoint of the two values in (9).  Therefore
\[
 \{|\Lambda_{2n}-q_n|<1/3\}\subseteq\{B_n=J_n\}.
\]
Equation~(12) would then give
\[
 P(B_n\ne J_n)\longrightarrow0,
\tag{14}
\]
contradicting the lower limit \(1/2\) in (8).  Thus no positive \(\mathcal F_{k_T}\)-measurable scale can satisfy the claimed suffix-QV stabilization for any deterministic \(k_T=o(T)\).  Since the argument imposed no upper or lower deterministic bound on \(\Lambda_T\), it applies in particular when \(\Lambda_T\) is uniformly bounded away from zero.

Finally, Proposition~\(\mathrm{prop:saco\text{-}theorem\text{-}5\text{-}14\text{-}counterexample}\) already verifies the variance-growth-only assumptions for this witness.  In its Saco specialization the scored set is \(\{1,\ldots,T\}\), so \(n_{\mathrm{eff}}=T\), and the ordered scored-set filtration in Definition~\(\mathrm{def:saco\text{-}early\text{-}stabilized\text{-}aipw\text{-}class}\) satisfies \(\mathcal G_{T,k}=\mathcal F_k\).  Hence \(k_T=o(n_{\mathrm{eff}})\) is exactly the regime just ruled out, and the witness cannot belong to that early-stabilized intersection class.